% !TEX root = actions.tex
\section{Patterns}
\label{patterns}

\def\ceq#1#2#3{\noindent\parbox[t]{22ex}{$\displaystyle #1$}\parbox{5ex}{\hfil $#2$}{$\displaystyle #3$}}

In the following $\U$ is some monster model of a complete theory $T$ of language $L$.
For simplicity $T$ is assumed to be one-sorted.

We define a new structure $\N$ whose domain is a subset of $\U^x$, where $x$ is some fixed tuple of variables.
The language of $\N$ is denoted by $L^{\rm\! p}$.
Though $L^{\rm\! p}$ depends on $x$, this will be left implicit.
We also use an auxiliary tuple $z=\langle z_i:i<\lambda\rangle$ of (so-called) parameter variables.

Fix a set of formulas $L_z\subseteq\Delta(x\,;z)\subseteq L_{x,z}$.
Let $V$ be some large set of variables that are copies of $x$. 
Finally, let $\Delta$ be the set of Boolean combinations of formulas in $\ \bigcup_{x'\in V}\Delta(x'\,;z)$.
The choice of $\Delta$ is justified in Remark~\ref{rem_first}.

Let $A\subseteq\U^x$ and $B\subseteq\U^z$.
We write $A,\Delta$ for the set of formulas $\phi(\bar a\,;z)$ for $\phi(\bar x\,;z)\in\Delta$ and $\bar a\in A^{\bar x}$, and we write $\Delta,B$ for the set of formulas $\phi(\bar x\,;b)$ for $\phi(\bar x\,;z)\in\Delta$ and $b\in B^z$.

Let $\bar a,\bar a'\in\big(\U^x\big)^n$ and  $b,b'\in\U^z$.
If $\phi(\bar a\,;b)\leftrightarrow\phi(\bar a'\,;b)$ holds for every $\phi(\bar x,z)\in\Delta$ and $b\in B^z$, we write \emph{$\bar a\equiv_{\Delta,B}\bar a'$.}
We define \emph{$b\equiv_{A,\Delta}b'$} similarly.
%We write $\Delta\mbox{-tp}(\bar a/M)$ for the set of formulas $\phi(\bar x\,;b)$ $\phi(\bar x,z)\in\Delta$ and $b\in M^z$.
% or equivalently $\Delta\mbox{-tp}(\bar a/M)=\Delta\mbox{-tp}(\bar a'/M)$ 

For every formula $\phi(x_0,\dots,x_{n-1}\,;z)\in\Delta$ the language $L^{\rm\! p}$ contains the $n$-ary relation \emph{$r_{\phi}$.}
Below by the symbol \emph{$L^{\rm\! p}$} we will also denote the set of \textit{atomic\/} formulas $r_{\phi}(x_0,\dots,x_{n-1})$ for any $x_0,\dots,x_{n-1}\in V$.
Note that equalities are \textit{not\/} included in $L^{\rm\! p}$.

Given pairs of models $M\preceq N$, where $N$ is $|M|^+$-saturated, we define a model $\N$ with domain $N^x$. 
% Eventually, we will proved that the main definition in this section, the definition of \textit{core,} does not depend on the choice of $M$.
% However, some of the intermediate notions introduced here do depend on $M$, though to avoid cluttering the notation we will omit reference to it.
We define the interpretation of $r_{\phi}$ in $\N$.
If $\bar a=a_0,\dots,a_{n-1}\in \N$, we declare $r_{\phi}(\bar a)$ true if $\phi(\bar a\,;b)$ holds for every $b\in M^z$ such that $\alpha(b)$.
We repeat the definition for emphasis

% \ceq{\hfill r_{\phi}(\bar a)}{\Leftrightarrow}{\forall z\in\alpha(M^z)\ \phi(\bar a\,;z).}
\ceq{\hfill r_{\phi}(\bar a)}{\Leftrightarrow}{\phi(\bar a\,;M^z)=M^z.}

The following is paramount.

\begin{remark}\label{rem_first}
  By the choice of $\Delta$, the truth of $r_{\phi}(\bar a)$ only depends on the equivalence class of $\equiv_{\Delta,M}$ of the \textit{components\/} of the tuple $\bar a$.
  Explicitely, if $\bar a=a_0,\dots,a_{n-1}$ and $\bar a'=a'_0,\dots,a'_{n-1}$ are such that $a_i\equiv_{\Delta,M} a'_i$ for every $i<n$, then \smallskip
  
  \ceq{\#\hfill r_{\phi}(\bar a)}{\leftrightarrow}{r_{\phi}(\bar a').} \smallskip

  % But we can stretch this further.
  % Write $a\equiv_{\{M\}}a'$ if $fa=a'$ for some $f\in{\rm Aut}(N/\{M\})$ i.e.\@ $a$ and $a'$ are conjugated by an automorphism that fixes $M$ setwise.
  % Then $\#$ holds whenever $a_i\equiv_{\{M\}} a'_i$ for every $i<n$.

\end{remark}

By the remark, it would be natural to define the interpretation of $L^{\rm\! p}$ on the quotient $\N/{\equiv_{\Delta,M}}$.
% or even on $\N/{\equiv_{\{M\}}}$~--~albeit the latter is more perplexing and will not be considered here.
But we prefer to keep the original domain $\N$ to have a common playground with $L$.

A \emph{p-type\/} is a set of formulas in $L^{\rm\! p}(\N)$.
We write \emph{$\textrm{p-tp}(\bar a)$} for the p-type of $\bar a$, i.e.\@ the set of formulas $r_{\phi}(\bar x)$ such that $r_{\phi}(\bar a)$ holds.

Note incidentally that the relation $x_0\equiv_{\Delta,M} x_1$ is definable by a p-type.
Namely the p-type containing the formulas $r_{\theta}(x_0,x_1)$ for every $\theta(x_0,x_1\,;z)$ of the form $\phi(x_0\,;z)\leftrightarrow\phi(x_1\,;z)$.



\begin{fact}\label{fact_p_qe}
  Formulas in $L^{\rm\! p}$ are closed under adjunction of dummy variables, disjunction, conjunction, and existential quantification.
\end{fact}

\begin{proof}
  The claim about dummy variables is clear, so we tacitatly use it below.
  
  We say that $\phi'(\bar x\,;z)$ is an \emph{alphabetic variant\/} of $\phi(\bar x\,;z)$ if it is obtained by replacing the vartiables effectively occurring in $\phi(\bar x\,;z)$ by other variables among those in $z$.
  Clearly, if $\phi'(\bar x\,;z)$ is an alphabetic variant of $\phi(\bar x\,;z)$, then $r_{\phi}(\bar x)$ and $r_{\phi'}(\bar x)$ are equivalent.

  Up to replacing $\phi(\bar x\,;z)$ with an alphabetic variant, we can assume that the $z$-variables that effectively occur in $\phi(\bar x\,;z)$ and are disjoint from those that occurin $\psi(\bar x\,;z)$.
  Therefore the following equivalences hold 
  
  \ceq{\hfill r_{\phi}(\bar a)\vee r_{\psi}(\bar a)}{\Leftrightarrow}{\forall z\in M^z\ \phi(\bar a\,;z)\ \ \vee\ \ \forall z\in M^z\ \psi(\bar a\,;z)}

  \ceq{}{\Leftrightarrow}{\forall z\in M^z\ \big[ \ \phi(\bar a\,;z)\, \vee\,\psi(\bar a\,;z)\big]}  

  \ceq{}{\Leftrightarrow}{r_{\phi\vee\psi}(\bar a)}

  A similar argument applies to the conjunction of two formulas.

  Finally, consider the formula $\exists y\,r_{\phi}(\bar x,y)$.
  Without loss of generality, we can assume that $\phi$ has the form

  \ceq{\hfill\phi(\bar x,y\,;z)}{=}{\bigvee_{i=1}^m\phi_i(\bar x,y\,;z)}

  where $\phi_i(\bar x,y\,;z)\,=\,\phi'_i(\bar x\,;z)\,\wedge\,\psi_i(y\,;z)$ and each $\phi'_i(\bar x\,;z)$ is a conjunction of formulas in $\bigcup_{x'\in V} L_{x'\,;z}$ or negation thereof.
  Closure under existential quantification follows from the following equivalences

  \ceq{\hfill\exists y\ r_{\phi}(\bar a, y)}{\Leftrightarrow}{\exists y\,\forall z\in M^z \bigvee_{i=1}^m\phi_i(\bar a,y\,;z)}

  \ceq{}{\Leftrightarrow}{\forall z\in M^z \bigvee_{i=1}^m\phi'_i(\bar a\,;z)\wedge\exists y\,\psi_i(y\,;z)\Big]}

  \ceq{}{\Leftrightarrow}{\vphantom{\bigvee_{i=1}^n}r_{\theta}(\bar a),}

  where 
  
  \ceq{\hfill\theta(\bar x\,;z)}{=}{\bigwedge_{i=1}^m\phi'_i(\bar x\,;z)\ \wedge\ \exists y\,\psi_i(y\,;z).}

  As $\exists y\,\psi_i(y\,;z)\in L_z\subseteq\Delta$, the above formula in $\Delta$ as required.
\end{proof}

\begin{fact}[ ($\infty$-p-saturation)]\label{fact_psat}
  Let $A\subseteq \N$ have arbitrary cardinality.
  Let $\bar x$ be a tuple of copies of $x$ of length $\le|\N|$.
  Then every p-type $p(\bar x)\subseteq L^{\rm\! p}(A)$ that is finitely consistent in $\N$ is realized in $\N$.
\end{fact}

\begin{proof}
  Let $A'\subseteq\N$ be a small set containing a representative of every $\equiv_{\Delta,M}$ equivalence class that intersects $A$.
  Replacing the elements of $A$ that occur in $p(\bar x)$ by elements of $A'$ we obtain an equivalent type $p'(\bar x)\subseteq L^{\rm\! p}(A')$.
  Note that every formula $r_{\phi}(\bar x)$ there is a $\Delta$-type $q(\bar x)\subseteq L(M)$ such that

  \ceq{\hfill \N\models r_{\phi}(\bar a)}{\Leftrightarrow}{N\models q(\bar a)}

  Therefore the saturation of $N$ implies that $p'(\bar x)$ is realized in $\N$.
\end{proof}

%%%%%%%%%%%%%%%%%%%%
%%%%%%%%%%%%%%%%%%%%
%%%%%%%%%%%%%%%%%%%%
%%%%%%%%%%%%%%%%%%%%
%%%%%%%%%%%%%%%%%%%%
%%%%%%%%%%%%%%%%%%%%
\section{Morphism}

A p-morphism is a map $f:\N\to\N$ such that for every $r_{\phi}(\bar x)\in L^{\rm p}$ and every $\bar a\in({\rm dom}f)^{\bar x}$

\ceq{\#\hfill\N\models r_{\phi}(\bar a)}{\Rightarrow}{\N\models r_{\phi}(f\bar a)}.

As $\equiv_{\Delta,M}$ is definable by a p-type, p-morphisms induce well-defined maps on the quotient $\N/{\equiv_{\Delta,M}}$.
We say that two p-morphisms $f$ and $g$ are \emph{equivalent,} and write $f\sim g$, if they induce the same map on $\N/{\equiv_{\Delta,M}}$.
We often confuse equivalent p-morphisms and we use the same symbol $f$ to denote the map induced on the quotient.

A total p-morphism is called a \emph{p-endomorphism.}
We write \emph{${\rm End}(\N)$} for the set of p-endomorphisms of $\N$.
We endow ${\rm End}(\N)$ with the topology that has as basic closed sets those of the form 

\ceq{\hfill\Big\{\ f\in{\rm End}(\N)}{:}{\N\ \models\ r_{\phi}\big(f\!a_1,\dots,f\!a_n,c_1,\dots,c_m\big)\Big\}}

for $\phi(x\,;z)\in\Delta$ and $a_1,\dots,a_n,c_1,\dots,c_m\in\N$.
Note that, by Remark~\ref{rem_first}, if $f\sim g$ then $f$ and $g$ belong to the same basic closed sets, therefore this topology naturally induces a topology on ${\rm End}(\N)/{\sim}$.

\begin{fact}
  With the above topology ${\rm End}(\N)$ is T$_1$ (i.e.\@ singletons are closed) and compact.
  Composition of p-morphisms is a left continuous operation.
\end{fact}

\begin{proof}
  We prove that singletons are closed.
  In fact

  \ceq{\hfill\big\{f\big\}}{=}{\bigcap_{a\in\N}\Big\{ g\in{\rm End}(\N)\ :\ a\equiv_{\Delta,M}ga\Big\}}

  therefore, as $\equiv_{\Delta,M}$ is defined by a p-type, $\{f\}$ is the intersection of closed sets.
  
  We prove compactness.
  Assume the following intersection is finitely consistent (we have slightly compressed the notation)

  % \ceq{\hfill\bigcap_{i<\lambda}\Big\{\ f\in{\rm End}(\N)}{:}{\N\ \models\ r_{\phi}\big(f\!a_1,\dots,f\!a_{n_i},a'_1,\dots,a'_m_{i}\big)\Big\}}


  \ceq{\#.\hfill}{~}{\bigcap_{i<\lambda}\Big\{\ f\in{\rm End}(\N)\ :\ \N\ \models\ r_i\big(f\!\bar a,\bar c\big)\Big\}.}

  then the p-type

  \ceq{\hfill\bar x\equiv_{\Delta, M}\bar a}{\cup}{\big\{r_i\big(\bar x,\bar c\big) \ :\ i<\lambda\big\}}

  is finitely consistent in $\N$.
  By $\infty$-p-saturation, it is realized by some $\bar a'$.
  By Fact~\ref{fact_phogeneity} there is a p-endomorphism $\N$ that maps $\bar a$ to $\bar a'$.
  This proves that the set in $\#$ is non empty.

  We prove left continuity of composition.
  That is for any given $f\in{\rm End}(\N)$ we prove that the function $\tau_f:{\rm End}(\N)\to{\rm End}(\N)$ that maps $g\mapsto g\circ f$ is continuous.
  Pick any basic closed set $\{h\, :\, r(h\bar a,\bar c)\}$.
  Then 
  
  \ceq{\hfill\tau_f^{-1}\big[\big\{h\,:\,r\big(h\bar a,\bar c\big)\big\}\big]}{=}{\big\{g\ :\ r\big(g(f\!\bar a),\bar c\big)\big\}}
  
  is closed, as required.
\end{proof}


% Though the switching between the two perspectives is intuitively obvious, it requires a fair amount of notation.

% Let $R\mathrel{\subseteq}\U^x\times\U^x$ be a binary relation.
% If $\bar a=\langle a_i:i<\lambda\rangle$ and $\bar b=\langle b_i:i<\lambda\rangle$ we write $\langle \bar a,\bar b\rangle\in R$ to abbreviate: $\langle a_i,b_i\rangle\in R$ for all $i<\lambda$.
% We say that $R$ is a \emph{p-morphism} if the following condition holds 

% \ceq{\#\hfill\U^x\models r_{\phi}(\bar a)}{\Rightarrow}{\U^x\models r_{\phi}(\bar b)}\hfill for every $r_{\phi}\in L^{\rm p}$ and every $\langle \bar a,\bar b\rangle\in R$.

% When the converse implication in \# also holds, i.e.\@ when also $R^{-1}$ is a p-morphism, we say that $R$ is a \emph{partial p-isomorphism.}
% When $\A\subseteq{\rm dom}\,R$, $\B\subseteq{\rm range}\,R$, and $R\cap (\A\times\B)$ is a partial p-isomorphism, we say that $R$ is an \emph{p-isomorphism} between $\A$ and $\B$ or, when $\A=\B$, a \emph{p-automorphism} of $\A$.

\begin{fact}\label{fact_phogeneity}
  Every p-morphism extends to a p-morphism defined on the whole of $\N$.
  Also, every p-isomorphism extends to a p-automorphism of $\N$.
\end{fact}

\begin{proof}
  Let $f:\N\to\N$ be a p-morphism.
  Let $\bar a$ be an enumeration of ${\rm dom} f$ and let $b\in\N$.
  Let $p(\bar x)=\textrm{p-tp}(\bar a)$ and let $q(\bar x,y)=\textrm{p-tp}(\bar a,b)$.
  By $\infty$-p-saturation $\exists y\,q(\bar x,y)$ is equivalent to the type containing the existential closure of the formulas in $q(\bar x,y)$.
  By Fact~\ref{fact_p_qe}, the latter is equivalent to a p-type, which is necessarily coicides with $p(\bar x)$.
  Then $q(f\bar a,y)$ is consistent hence by $\infty$-p-saturation it is realized by some $c\in\N$.
  Then $f\cup\{\langle b,c\rangle\}$ is a p-morphism defined in $b$ that extends $f$.

  The second claim is obtained by back-and-forth.
\end{proof}

%%%%%%%%%%%%%%%%%
%%%%%%%%%%%%%%%%%
%%%%%%%%%%%%%%%%%
%%%%%%%%%%%%%%%%%
%%%%%%%%%%%%%%%%%
%%%%%%%%%%%%%%%%%
\section{Duality}

In this section $\hat m$ is an enumeration of $M$ and $|z|=|\hat m|$.
For a more convenient notation, $x$ is taken of infinite length~--~then, up to equivalence, formulas in $L_z(\U)$ have the form $\phi(a\,;z)$ for some $\phi(x\,;z)\in L$ and $a\in\U^x$.

\begin{remark}
  We stress a simple fact which is used in the proof below. 
  Recall that $\phi'(a\,;z)$ is an alphabetic variant of $\phi(a\,;z)$ if it is obtained by replacing the vartiables effectively occurring in $\phi(a\,;z)$ by other variables among those in $z$.
  The following are clearly equivalent
  \begin{itemize}
    \item[1.] $\phi(a\,;M^z)=M^z$
    \item[2.] $\phi'(a\,;\hat m)$ holds for every alphabetic variant of $\phi(a\,;z)$.
  \end{itemize}
\end{remark}

Let $\B=\big\{b\in\U^z\ :\ \hat m\equiv b\nonfork_MN\big\}$, where $\equiv$ is the elementary equivalence in $L$ however, as $L_z\subseteq\Delta$, we could replace it with $\equiv_{\varnothing,\Delta}$.
The $N,\Delta$-topology on $\B$ is the one generated by the closed sets of the form $\{\ b\in\B\ :\ \phi(\bar a\,;b)\ \}$ for $\phi(\bar x\,;z)\in\Delta$ and $\bar a\in\N^{\bar x}$.
The $N,\Delta$-topology on $\B/{\equiv_{N,\Delta}}$ is the quotient of that on $\B$.


\begin{fact}
  The topological spaces ${\rm End}(\N)/{\sim}$ and $\B/{\equiv_{N,\Delta}}$ are homeomorphic.
\end{fact}

\begin{proof}
  In 1 below, we define a map $f\mapsto b_f$ from ${\rm End}(\N)$ into $\B$. In 2, we define a map $b\mapsto f_b$ from $\B$ into ${\rm End}(\N)$.
  The reader may check that these maps induce maps between ${\rm End}(\N)/{\sim}$ and $\B/{\equiv_{N,\Delta}}$ that are each other's inverse.
  In 3 we prove the continuity of the second map.

  1.\ 
  Let $f: \N\to \N$ be an p-epimorphism.
  Let $p(z)$ be the type containing the formulas $\phi(a\,;z)$ for $\phi(x\,;z)\in\Delta$ and $a\in \N$ such that $\phi(f\!a\,;\hat m)$.
  Clearly $p(z)\vdash\hat m\equiv z$ and is complete over $N,\Delta$.
  We prove that $p(z)$ is finitely satisfied in $M$, then $b_f$ is some/any realization of $p(z)$.

  If for a contradiction that $p(z)$ is not finitely satisfied in $M$, then $\psi(a_1,\dots,a_n\,;M^z)=\varnothing$, where 
  
  \ceq{\hfill\psi(x_1,\dots,x_n\,;z)}{=}{\bigwedge_{i=n}^n\phi_i(x_i\,;z)}
  
  for some $\phi_i(x\,;z)\in L$ and $a_i\in \N$ such that $\phi_i(f\!a_i\,;\hat m)$ for every $i$.
  Then $r_{\neg\psi}(a_1,\dots,a_n)$ and $\neg r_{\neg\psi}(f\!a_1,\dots,f\!a_n)$.
  This is not possible because $f$ is a p-morphism.

  2.\ 
  Let $b\in\B$ be given.
  We define 
  
  \ceq{\hfill p(x,y,z)}{=}{\big\{\phi(y\,;z)\rightarrow\phi(x\,;\hat m)\ :\ \phi(x\,;z)\in\Delta\big\}.}

  Note that $p(x,a,b)$ is a p-type, for any $a\in\N$ and $b\in\B$.
  We claim that if $a\equiv_{\Delta,M}a'$ then $p(x,a,b)=p(x,a',b)$.
  In fact, $\phi(a\,;b)\nleftrightarrow\phi(a'\,;b)$ would contradict $b\nonfork_MN$.
  Also, note that $p(x,a,b)$ is consistent, otherwise $\neg\exists x\bigwedge_i\phi_i(x\,;\hat m)$, for some finitely many formulas $\phi_i(x\,;z)\in\Delta$ such that $\phi_i(a\,;b)$.
  But this would contradict $b\equiv\hat m$.

  Then we can define $f_b(a)$ to be some/any realization of $p(x,a,b)$.

  We prove that $f_b$ is a p-morphism.
  Let $\psi(x_1,\dots,x_n\,;z)\in\Delta$ and $a_1,\dots,a_n\in\U^x$ be such that $r_{\psi}(a_1,\dots,a_n)$.
  % We can assume $\alpha(M^z)\neq\varnothing$, otherwise $r_{\psi}(f\!a_1,\dots,f\!a_n)$ trivially holds.
  Assume

  \ceq{\hfill\psi(x_1,\dots,x_n\,;z)}{=}{\bigvee_{h=1}^m\bigwedge_{k=1}^m\phi_{h,k}(x_{i_{h,k}}\,;z)}

  for some $i_{h,k}\in\{1,\dots,n\}$.
  After re-indexing, from $r_{\psi}(a_1,\dots,a_n)$ we obtain $\phi_k(a_k;M^z)=M^z$ for $k=1,\dots,m$.
  Then from $b\nonfork_MN$ we obtain that $b$ satiesfies $\phi_k(a_k\,;z)$ and all its alphabetic variants.
  Therefore $\hat m$ satisfies $\phi_k(f_ba_k\,;z)$ and all its alphabetic variants.
  Therefore $\phi_k(f_ba_k;M^z)=M^z$, and $r_{\psi}(f_ba_1,\dots,f_ba_n)$ follows.

  % 3.\ We prove that the map defined in 1 is continuous.
  % Pick any basic closed set in $\B$, say $\big\{b\,:\,\phi(\bar a\,;b)\big\}$.
  % We claim that the set $\big\{f\,:\,\phi(\bar a\,;b_f)\big\}$ is closed in ${\rm End}(\N)$. Indeed, by the definition of $b_f$ in 1

  % \ceq{\hfill\big\{f\ :\ \phi\big(\bar a\,;b_f\big)\}}{=}{\Big\{f\ :\ \phi\big(f\bar a\,;\hat m\big)\Big\}}

  % \ceq{\#}{=}{\Big\{f\ :\ r_{\psi}\big(f\bar a,c\big)\Big\}}

  % where $\psi(\bar x,y;z)$ and $c$ are choosen so that $r_{\psi}(\bar x,c)$ is equivalent to $\phi(\bar x\,;\hat m)$.
  % As the set in $\#$ is closed in ${\rm End}(\N)$, continuity of the map $f\mapsto b_f$ follows.

  % We show how to choose $\psi(\bar x,y;z)$ and $c$.
  % Let $y$ be a copy of $x$.
  % The formula $\psi(\bar x,y;z)$ is obtained from $\phi(\bar x\,;z)$ by replacing the variables in the tuple $z$ effectively occurring in $\phi(\bar x\,;z)$ with variables in the tuple $y$.
  % Note that $z$ is dummy in $\psi(\bar x,y\,;z)$.
  % Finally, we choose $c\in M^{y}$ such that $\psi(\bar x,c\,;z)$ is equivalent to $\phi(\bar x\,;\hat m)$.

  3.\ 
  We prove that the map defined in 2 is continuous.
  Pick any basic closed set in ${\rm End}(\N)$, say $\big\{ h\, :\, r(h\bar a,\bar c)\big\}$.
  We claim that the set $\big\{b\, :\, r(f_b\bar a,\bar c\big)\}$ is closed. Indeed, by the definition of $f_b$ in 2,

  \ceq{\hfill\big\{b\ :\ r\big(f_b\bar a,\bar c\big)\big\}}{=}{\Big\{b\ :\ b\models \exists\bar x\ \Big[\bigwedge_{i=1}^np(x_i,a_i,b)\wedge r\big(\bar x,\bar c\big)\Big]\Big\}}

  is closed in $\B$ because $r\big(\bar x,\bar c\big)$ is definable by a $N,\Delta$-type.
\end{proof}

% We define a semigroup operation on $\B$.
% Namely, if $f_b,f_c:\to$ are the p-morphisms the correspond to $b,c\in\B$ we write $b{*}c$ for the element of $B$ that corresponds to the p-morphism $f_b\circ f_c$.

% On $\B/{\equiv_{\Delta,M}}$ the operation $*$ is well-defined and turns it into a semigroup.



% Let $p(z)$ and $q(z)$ be the types that correspond to the p-morphisms $f_p$, respectively $f_q$.
% Pick any sufficiently saturated model $N\succeq M$ and $c\models q(z)\restriction N$.
% Then the type $(p\otimes q)(z,c)$ corresponds to the p-morphisms $f_p\circ f_q$.

% 

%

% $(p\otimes q)(z,c)=\{\phi(z,c)\ \ :\ \ \phi(z,c')\in p \textrm{ for some/any } c'\models q(z)\restriction N\}$

% Then $\phi(a',c)$ holds for every $a'\in f_q^{-1}[N]$.

% $z'$ be a copy of $z$
% Say $f,g: S(M)\to S(M)$ and $h=gf$.
% Let $m,n$ enumerate $M$.
% Write $p(z)=\Delta\mbox{-tp}(m/\varnothing)$.
% By Lemma 3.14, we have that $g$ and $f$ correspond to global extensions $p_f$ and $p_g$. 
% Explicitely:

% $p_f=\{\phi(a,z)\ :\ φ(fa,m),\ a\in\U^x\}$

% $p_g=\{\phi(a,z)\ :\ φ(ga,m),\ a\in\U^x\} = \{\phi(fa,z)\ :\ φ(gfa,m),\ a\in\U^x\}$

% $p_{gf}=\{\phi(a,z)\ :\ φ(gfa,m),\ a\in\U^x\}$

% % f(tp(e/A))={φ(x,b) : φ(e',y)∈f*p for some/any e' ≡_A e}={φ(x,b) : φ(e,y)∈f*p}
% % g*q={ φ(e,z) : φ(x,c)∈ g(tp(e/B)) }
% % g(tp(e/B))={φ(x,c) : φ(e',z)∈g*q for some/any e' ≡_B e}={φ(x,c) : φ(e,z)∈g*q}

% Now, h corresponds to a global extension h*q of q. Let's write h*q in terms of f*p and g*q:
% h*q={φ(e,z) : φ(e,z)\in gf(tp(e/A))}={φ(e,z) : φ(e',z)∈g*q for some/any e' realising f(tp(e/A))}={φ(e,z) : φ(e',z)∈g*q for some/any e' such that ⊧θ(e',b) for any θ(e,y)∈f*p}={φ(e,z) : φ(e',z)∈g*q for some/any e' such that e'b≡eb' for any b' realising f*p|e}

% Hence, for p(y) global A-invariant type extending tp(b/ø) and q(z) global B-invariant type extending tp(c/ø), define
% q⊗p={φ(e,z) : φ(e',z)∈q for some/any e' such that e'b≡eb' for any b' realising p|e}.
% We can further check that q⊗p is finitely satisfiable in A if q is finitely satisfiable in B and p is finitely satisfiable in A.  

%%%%%%%%%%%%%%%%%%%%
%%%%%%%%%%%%%%%%%%%%
%%%%%%%%%%%%%%%%%%%%
%%%%%%%%%%%%%%%%%%%%
%%%%%%%%%%%%%%%%%%%%
\section{Semigroups}

Assume $\U$ is a unital semigroup (a monoid) w.r.t.\@ some definable operation which we denote by $\cdot$.
The unit is denoted by $1$.
In this section we could work with $|x|=|z|=1$ but, for uniformity with the previous section we take $|x|=|z|=\omega$.
We write $z\cdot x$ for coordinatewise multiplication.
In this section $\Delta(x,z)$ contains the formulas $\phi(z\cdot x)$ for all $\phi(x)\in L$.

\begin{fact}
  For $b\in\B$, let $g_b:\N\to\N$ be the function that maps $a\mapsto b\cdot a$.
  Then $b\mapsto g_b$ is an homeomorphism between $\B/{\equiv_{\Delta,M}}$ and ${\rm End}(\N)/{\sim}$.
\end{fact}

\begin{proof}
  We prove that $g_b$ is a p-morphism.
  Assume $r_{\psi,\alpha}(a_1,\dots,a_n)$.
  We prove $r_{\psi,\alpha}(g_ba_1,\dots,g_ba_n)$.
  Assume that for some $i_{h,k}\in\{1,\dots,n\}$

  \ceq{\hfill\psi(x_1,\dots,x_n\,;z)}{=}{\bigvee_{h=1}^m\bigwedge_{k=1}^m\phi_{h,k}(\,z\cdot\, x_{i_{h,k}})}

  for some $\phi_{h,k}(z\cdot x)\in\Delta$.
  From $r_{\psi}(a_1,\dots,a_n)$, after re-indexing, we obtain that $\phi_k(M^z\cdot a_k)=M^z$ for $k=1,\dots,m$.
  As $b\nonfork_MN$, we immediately obtain $\phi_k(M^z\cdot b\cdot a_k)=M^z$ for all $k$, so $r_{\psi,\alpha}(g_ba_1,\dots,g_ba_n)$ follows.
\end{proof}

%%%%%%%%%%%%%%%%%%%%
%%%%%%%%%%%%%%%%%%%%
%%%%%%%%%%%%%%%%%%%%
%%%%%%%%%%%%%%%%%%%%
%%%%%%%%%%%%%%%%%%%%
\section{Cores}

% \begin{fact}\label{fact_reduced_funct} 
%   Let $R$ be a p-morphism and let $f\subseteq R$ be a reduced p-morphism.
%   Then there is a reduced p-morphism $f'$ extendiung $f$ and such that 

%   \ceq{\hfill\big(\mathrel{\equiv_{\Delta,M}}\circ \mathrel{R}\circ\mathrel{\equiv_{\Delta,M}}\big)}{=}{\big(\mathrel{\equiv_{\Delta,M}}\circ \mathrel{f'}\circ\mathrel{\equiv_{\Delta,M}}\big).}
% \end{fact}

% \begin{proof}
%   Let $R$ be a p-morphism.
%   Pick a reduced p-morphism $f$ such that $R\ \subseteq\ \mathrel{\equiv_{\Delta,M}}\circ \mathrel{f}\circ\mathrel{\equiv_{\Delta,M}}$.
%   Let $\bar a$ be an enumeration of ${\rm dom} f$ and let $b\in\U^x$.
%   Let $p(\bar x)=\textrm{p-tp}(\bar a)$ and let $q(\bar x,y)=\textrm{p-tp}(\bar a,b)$.
%   By $\infty$-p-saturation and Fact~\ref{fact_p_qe}, $p(\bar x)$ is equivalent to $\exists y\,q(\bar x,y)$.
%   Then $q(f\bar a,y)$ is consistent so $f$ can be extended to a p-morphism defined on $b$.
%   The second claim follows by back-and-forth.
% \end{proof}

A set $A\subseteq\N$ is \emph{p-maximal} if it satisfies the following (obviously) equivalent conditions for every $\bar a\in A^{<\omega}$ and every p-morphism $f$ defined on $\bar a$:

\begin{itemize}
  \item[a.] $r_{\phi}(\,\bar a\,)\ \ \Leftrightarrow\ \ r_{\phi}(\,f\!\bar a\;)$ \hfill
  for every $r_{\phi}\in L^{\rm\! p}$

  \item[b.] $\textrm{p-tp}(\,\bar a\,)\ \ =\ \ \textrm{p-tp}(\,f\!\bar a\;)$.
\end{itemize}

% By Zorn's lemma, every set satisfying the above conditions is contained in a p-maximal set.

\begin{fact}\label{fact_pmaximal}
  The following are equivalent
  \begin{itemize}
    \item[1.] $A\subseteq\N$ is p-maximal
    \item[2.] $p(\bar x)=\textrm{p-tp}(\bar a)$ is a maximally consistent p-type for every $\bar a\in A^{\bar x}$ and every finite $\bar x$
    \item[3.] $p(\bar y)=\textrm{p-tp}(\bar c)$ is a maximally consistent p-type for $\bar c$ any enumeration of $A$
    \item[4.] every p-morphism defined on $A$ is a p-isomorphism between $A$ and its image.
  \end{itemize}
\end{fact}

\begin{proof}
  1$\Rightarrow$2. Negate 2. Suppose $\textrm{p-tp}(\bar a)$ is not maximal for some finite $\bar a\in A^{\bar x}$.
  Pick some $\bar a'\in A^{\bar x}$ such that $\textrm{p-tp}(\bar a)\subset \textrm{p-tp}(\bar a')$.
  Then $f=\{\langle \bar a,\bar a'\rangle\}$ is a p-morphism that contradicts the p-maximality of $A$.
  
  2$\Rightarrow$3. If $\bar c$ is an enumeration of $A$ that witnesses the failure of 3, then some finite subtuple of $\bar c$ witnesses the failure of 2.

  3$\Rightarrow$4. Negate 4. Let $f$ be a p-morphism defined on $A$ that is not a p-isomorphism.
  Let $\bar a\in A^{\bar x}$ be such that  $\textrm{p-tp}(\bar a)\smallsetminus \textrm{p-tp}(f\bar a)\neq\varnothing$.
  To contradict 3, it suffices to extend $f$ to a p-morphism defined on $A$.
  This is possible by $\infty$-p-homogeneity.

  4$\Rightarrow$1. Clear.
\end{proof}

\begin{example}
  Assume that $a\in\N$ is such that $\phi(a\,;M^z)$ is definable over ${\rm acl^{eq}}\varnothing$ for every $\phi(x\,;z)\in L$.
  Then the orbit of $a$ under ${\rm Aut}(\U/\{M\})$ is p-maximal.
\end{example}

\begin{proof}
  First consider the case when $\phi(a\,;M^z)$ is definable over $\varnothing$.
  We prove that $\{a\}$ is p-maximal.

  Let $\phi'(z)\in L$ be such that $\phi(a\,;M^z)=\phi'(M^z)$.
  Then, if we set $\phi''(x\,;z)=\phi(x\,;z)\leftrightarrow\phi'(z)$, the formula $r_{\phi''}(a)$ holds.

  If $f$ is a p-morphism then $r_{\phi''}(fa)$ also holds.
  Then $a\equiv_{\Delta,M}f\!a$ follows.
  This shows that p-morphism are trivial, i.e.\@ equivalent to the identity, and therefore p-isomorphisms.

  Now assume that $\phi(a\,;M^z)$ is definable over ${\rm acl^{eq}}\varnothing$.???
\end{proof}

A set $\J\subseteq\N$ is called a \emph{core} if it is p-maximal and, among the p-maximal sets, it is maximal w.r.t.\@ inclusion.

\begin{fact}\label{fact_onto}
  Let $\J$ be a core.
  Let $\A$ be a p-maximal set.
  Then every p-morphism $f:\A\to\J$ extends to a p-isomorphism between $\A$ and $\J$. 
\end{fact}

\begin{proof}
  Reason as in the proof of Fact~\ref{fact_phogeneity} with the same notation but choose $b\in\A$.
  Let $f'$ be a p-morphism defined on $\A$ that extends $f\cup\{\langle b,c\rangle\}$.
  Then $f'[\A]$, i.e.\@ the image of $\A$ under $f'$, is a p-maximal set that contains $\J$.
  By maximality $f'[\A]$ coincides with $\J$, hence $c\in\J$.
\end{proof}

By back-and-forth we obtain the following.

\begin{corollary}[ ($\infty$-p-homogeneity of the core)]\label{fact_corephogeneity}
  If $\J$ and $\J'$ are cores, then every p-morphism $f:\J\to\J'$ extends to a p-isomorphism between $\J$ and $\J'$.
\end{corollary}

A \emph{p-retraction} of $\N$ onto $\A$ is a total p-morphism $f:\N\to\A$ that is equivalent to the identity on $\A$.

\begin{theorem}
  If $\J$ and is a core then there is a p-retraction of $\N$ onto $\J$
\end{theorem}

\begin{proof}
  Let $\bar c$ enumerate $\N$.
  Let $\bar a$ be a tuple such that $\textrm{p-tp}(\bar a)\supseteq \textrm{p-tp}(\bar c)$ is p-maximal.
  Let $f\ =\{\langle\bar c,\bar a\rangle\}$.
  Then $f:\N\to\A$ is a p-morphism onto the p-maximal set $\A={\rm range(\bar a)}$.
  As $\A$ is  p-maximal, by Fact~\ref{fact_onto} we can extend $f^{-1}$ to a p-isomorphism $f'$ defined on $\A$.
  It is easy to verify that $f'\circ f$ is a retraction.
\end{proof}

\begin{fact}
  The following are equivalent
  \begin{itemize}
    \item[1.] $a\in\bigcap\big\{\,\J\ :\ \J \textrm{ is a core}\big\}$
    \item[2.] $\phi(a\,;M^z)$ is definable over ${\rm acl^{eq}}\,\varnothing$, for every formula $\phi(x\,;z)\in L$.
  \end{itemize} 
\end{fact}

\begin{proof}
  ???
\end{proof}

\begin{fact}
  $\J$ is compact subset of $\N$.
\end{fact}

\begin{proof}
  ???
\end{proof}

\section{Esperimenti}

Let $M$ be a saturated model and let $G={\rm Aut}(M)$.

The group $G$ acts on $\U^x/{\equiv_M}$ in the natural way. 


